\documentclass[10pt,a4paper]{amsart}
\usepackage[margin=19mm]{geometry}
\usepackage{amsmath,amssymb,mathtools}
\usepackage[scr=rsfs,cal=euler]{mathalfa}
\usepackage{xfrac}
\usepackage[unicode,hidelinks,bookmarks=false]{hyperref}
\newcommand{\ie}{i.e.\ }
\newcommand{\cf}{cf.\ }
\newcommand{\resp}{respectively\ }
\newcommand{\Subsection}{\subsection}
\providecommand{\nobreakdash}{\nobreak\hbox{-}}
\setlength{\emergencystretch}{2em}
\multlinegap0pt
\newtheorem{theorem}{Theorem}[section]
\newtheorem{proposition}[theorem]{Proposition}
\newtheorem{corollary}[theorem]{Corollary}
\newtheorem{lemma}[theorem]{Lemma}
\newtheorem{remark}[theorem]{Remark}
\newtheorem{problem}[theorem]{Problem}
\newtheorem{conjecture}[theorem]{Conjecture}
\newtheorem{question}{Question}
\newtheorem{claim}{Claim}
\newtheorem{observation}{Observation}
\theoremstyle{definition}
\newtheorem{definition}[theorem]{Definition}
\title{Symmetric polytopes whose automorphism groups are 2-groups: finite regular construction}
\author{Gabe Cunningham, Yan-Quan Feng, Dong-Dong Hou, Egon Schulte}
\date{Source: arXiv:2512.15511v1; CC BY 4.0}
\begin{document}\maketitle
\noindent\textit{Counting scope.} The sole counted claim is the original theorem \texttt{flat44thm} (native source lines 354--362), with its complete author proof (364--419). Original definitions (100--135 and 148), the complete source-local FAP inheritance argument and its elementary supporting lemma (156--301), and toroidal input definitions and proof (306--349) are uncounted support. The mix theorem is expressly cited by the authors to McMullen--Schulte, Theorem 4F9; its proof is a stated prerequisite, not newly proved here. All retained mathematical passages are copied without replacements or omissions. Equation and theorem numbers follow this standalone wrapper; original labels and citation keys retain their targets. No figures occur in these spans. The title block credits all four source authors, including Egon Schulte whose original author block begins with a line break.\par\medskip
\noindent\textit{Source and licence.} The fixed author source is arXiv:2512.15511v1. The package carries the version-specific CC BY 4.0 notice and original bibliography entries; no publisher class, upstream script, font or external artwork is redistributed.\par\medskip
\section{Original definitions and support}
In this section we briefly review basic concepts and results about abstract regular polytopes, mostly following \cite[Chs. 2,\,3]{McMSch2002}. 

An (\emph{abstract\/}) \textit{polytope of rank\/} $d$ is a partially ordered set $\mathcal{P}$ with a strictly monotone rank function that takes values in $\{-1,0, \ldots, d\}$. An element of rank~$j$ is called a \textit{$j$-face\/} of $\mathcal{P}$, and a face of rank $0$, $1$ or $d-1$ is also called a \textit{vertex\/}, \textit{edge\/} or \textit{facet\/}, respectively. The polytope $\mathcal{P}$ has a unique least face $F_{-1}$ (of rank $-1$) and a unique greatest face $F_d$ (of rank $d$), called \textit{improper\/} faces of $\mathcal{P}$. A \textit{flag\/} is a maximal totally ordered subset of $\mathcal{P}$. Each flag contains exactly $d+2$ faces, one for each rank, including $F_{-1}$ and $F_{d}$. Two flags are said to be \textit{adjacent} if they differ in one face; if this face is a $j$-face, we also call the flags $j$-adjacent. A polytope $\mathcal{P}$ is \textit{strongly flag-connected}, in the sense that any two flags $\Phi$ and $\Psi$ of $\mathcal{P}$ can be joined by a finite sequence of successively adjacent flags, each containing $\Phi \cap \Psi$. Finally, $\mathcal{P}$ satisfies the \textit{diamond condition\/}:\ whenever $F \leq G$, with $F$ a $(j-1)$-face and $G$ a $(j+1)$-face for some~$j$, there are exactly two $j$-faces $H$ with $F \leq H \leq G$. The diamond condition, rephrased, is saying that every flag $\Phi$ has a unique $j$-adjacent flag for each $j=0, \dots, d-1$, denoted $\Phi^j$.  An abstract polytope of rank 3 is also called an (\emph{abstract}) \emph{polyhedron}.

For any two faces $F$ of rank $j$ and $G$ of rank $k$ with $F \leq G$, we call
\[G/F := \{ H \in \mathcal{P}\, | \, F \leq H \leq G \}\] 
a \textit{section} of $\mathcal{P}$. This is a $(k-j-1)$-polytope in its own right. We often identify a face $F$ with the section $F/F_{-1}$. For a $j$-face $F$ the section $F_{d}/F$ is called the \textit{co-$j$-face\/} (or just \textit{co-face\/}) of $\mathcal{P}$ at $F$, or the \emph{vertex-figure\/} or \emph{edge-figure\/} of $\mathcal{P}$ at $F$ if $F$ is a vertex or an edge, respectively. A co-$j$-face has rank $d-j-1$.

A polytope $\mathcal{P}$ is said to be \emph{regular\/} if its automorphism group $\Gamma(\mathcal{P})$ is transitive on the flags. In this case $\Gamma(\mathcal{P})$ acts regularly on the flags, as the flag stabilizers in the automorphism group of a polytope are necessarily trivial. If $\mathcal{P}$ is a regular polytope, then $\Gamma(\mathcal{P})$ acts transitively on the faces of each rank, all sections of $\mathcal{P}$ are also regular, and comparable sections (defined by pairs of faces of matching ranks) are isomorphic. A regular polytope is said to be of (Schl\"afli) type $\{p_1,\ldots,p_{d-1}\}$ if for each $i=1,\ldots,d-1$ and each incident pair $F,G$ of $(i-2)$-face and $(i+1)$-face, the section $G/F$ is isomorphic to a $p_i$-gon (possibly $p_{i}=\infty$). 

The group $\Gamma(\mathcal{P})$ of a regular polytope $\mathcal{P}$ has a well-behaved system of \textit{distinguished generators\/}. In fact, $\Gamma(\mathcal{P})$ is generated by involutions $\rho_0,\ldots,\rho_{d-1}$, where $\rho_i$ maps a (fixed) \textit{base\/} flag $\Phi$ to its $i$-adjacent flag $\Phi^i$. These generators satisfy (at least) the standard Coxeter-type relations for Coxeter groups with string diagrams, 
\begin{equation}
\label{standardrel}
(\rho_i \rho_j)^{p_{ij}} = 1 \;\; \textrm{ for } i,j=0, \ldots,d-1,
\end{equation}
where $p_{ii}=1$, $p_{ji} = p_{ij} =: p_{i+1}$ if $j=i+1$, and $p_{ij}=2$ otherwise. The numbers $p_j$ are the entries of the Schl\"afli symbol $\{p_{1},\ldots,p_{d-1}\}$. Moreover, the following {\em intersection property\/} holds,
\begin{equation}
\label{intprop}
\langle \rho_i \mid i \in K \rangle \cap \langle \rho_i \mid i \in J \rangle
= \langle \rho_i \mid i \in {K \cap J} \rangle
  \;\; \textrm{ for } K,J \subseteq \{0,1,\ldots,d-1\}.
\end{equation}
Further, the elements $\sigma_{1},\ldots,\sigma_{d-1}$, with $\sigma_{i}:=\rho_{i-1}\rho_{i}$ for $i=1,\ldots,d-1$, generate the {\em rotation subgroup\/} $\Gamma^{+}(\mathcal{P})$ of $\Gamma(\mathcal{P})$, which is of index at most~$2$. This subgroup consists of all elements of $\Gamma(\mathcal{P})$ that can be written as a product of an even number of generators $\rho_i$. A regular polytope $\mathcal{P}$ is {\em directly regular\/} if the index of $\Gamma^{+}(\mathcal{P})$ in $\Gamma(\mathcal{P})$ is $2$.

A \textit{string C-group\/} is a group $\Gamma = \langle \rho_0,\ldots, \rho_{d-1} \rangle$ whose generators satisfy (\ref{standardrel}) and the intersection property (\ref{intprop}); here, the ``C'' stands for ``Coxeter'', though not every string C-group is a Coxeter group. The string C-groups are precisely the automorphism groups of regular polytopes (see \cite[Ch. 2E]{McMSch2002}). We usually identify a regular polytope with its automorphism (string C-) group.

We write $[p_{1},p_{2},\ldots,p_{d-1}]$ for the Coxeter group whose underlying Coxeter diagram is a string with $d$ nodes and $d-1$ branches labeled $p_{1},p_{2},\ldots,p_{d-1}$. (Here we regard the $i$-th branch of the string as missing if $p_{i}=2$.) This is a string C-group and is the automorphism group of the \textit{universal} regular $n$-polytope $\{p_{1},\ldots,p_{d-1}\}$ (see \cite[Ch. 3D]{McMSch2002}). 

Assembling regular polytopes from polytopes of lower rank is usually a difficult problem. If a regular $d$-polytope $\mathcal{P}$ has facets isomorphic to $\mathcal{P}_1$ and vertex-figures isomorphic to $\mathcal{P}_2$, then $\mathcal{P}_1$ and $\mathcal{P}_2$ must necessarily be $(d-1)$-polytopes that are \textit{compatible}, in the sense that the vertex-figures of $\mathcal{P}_1$ must be isomorphic to the facets of $\mathcal{P}_2$. The compatibility of a pair of regular $(d-1)$-polytopes $\mathcal{P}_1$ and $\mathcal{P}_2$ is only necessary but not sufficient to guarantee the existence of a regular $d$-polytope with facets $\mathcal{P}_1$ and vertex-figures $\mathcal{P}_2$. However, if such a polytope indeed exists, then there is also a universal such polytope, denoted $\{\mathcal{P}_1,\mathcal{P}_2\}$, which covers all other regular polytopes with facets $\mathcal{P}_1$ and vertex-figures $\mathcal{P}_2$ (see \cite[Thm.\,4A2]{McMSch2002}). Generalizing the concept of the Schl\"afli symbol, we slightly abuse notation and say that a regular polytope $\mathcal{P}$ is of \textit{type\/} $\{\mathcal{P}_1,\mathcal{P}_2\}$ if $\mathcal{P}$ has facets $\mathcal{P}_1$ and vertex-figures $\mathcal{P}_2$.

We will also use a \textit{generalized Schl\"afli symbol\/} recording the structure of the sections of rank 3. Recall that if $\Phi=\{F_{-1},F_0,\ldots,F_d\}$ is a flag of a regular $d$-polytope $\mathcal P$, the sequence of sections of rank 3 determined by $\Phi$, 
\begin{equation}
\label{rk3sections}
F_3/F_{-1},\ F_4/F_0,\ F_5/F_1,\ \ldots,\ F_{d}/F_{d-4},
\end{equation}
captures the local geometry of the sections of rank 3 of $\mathcal P$, in the sense that each section $G/F$ given by a $j$-face $G$ and a $(j-4)$-face $F$ with $F\leq G$ is isomorphic to $F_j/F_{j-4}$. We then say that $\mathcal{P}$ is of type $\{\mathcal{P}_3,\ldots,\mathcal{P}_{d}\}$ if $\mathcal{P}_j$ is a regular 3-polytope isomorphic to the section $F_{j}/F_{j-4}$ of $\mathcal{P}$ for each $j=3,\ldots,d$. If $d=4$, this is consistent with our previous notation (apart from renaming the subscripts). Note that our notation is not meant to imply universality of $\mathcal{P}$ among regular polytopes with the same generalized Schl\"afli symbol.

A $d$-polytope $\mathcal{P}$ is called {\em flat\/} if each of its vertices is incident with each of its facets.  More generally, if $0 \leq k < l \leq d-1$, then $\mathcal{P}$ is said to be {\em $(k,l)$-flat\/} if each of its $k$-faces is incident with each of its $l$-faces. A flat polytope is $(0,d-1)$-flat.

In this section, we briefly review the flat amalgamation property, or FAP for short, as well as the amalgamation process for regular polytopes described in \cite[Sects.\,4E,\,4F]{McMSch2002}.

Let $\Gamma = \langle\rho_{0},\ldots,\rho_{d-1}\rangle$ be any group generated by involutions such that $(\rho_{i}\rho_{j})^{2} = 1$ for all~$i,j$ with $|i-j| \geq 2$. We usually take $\Gamma$ to be the group of a regular polytope.  Note that if $\varphi\in \Gamma$ is such that 
$\varphi = \varphi_{0}\rho_{j_{1}}\varphi_{1}\rho_{j_{2}} \cdots 
\varphi_{m-1}\rho_{j_{m}}\varphi_{m}$, 
then we can write 
\[\varphi  = \psi_{0}(\psi_{1}^{-1}\rho_{j_{1}}\psi_{1}) \cdots
(\psi_{m}^{-1}\rho_{j_{m}}\psi_{m}),  \]
where  $\psi_{i} := \varphi_{i}\varphi_{i+1} \cdots \varphi_{m}$ for $i = 0,\ldots,m$. 
Here, if $\varphi_{0},\ldots,\varphi_{m}$ belong to some specified subgroup
of $\Gamma$, then $\psi_{0},\psi_{1},\ldots,\psi_{m}$ also belong to that subgroup.

For $k = 0,\ldots,d-1$ define
\[\begin{array}{rclcl}
N_{k}^{+} & := & N_{k,\ldots,d-1} & := & \langle\varphi^{-1}\rho_{i}\varphi \mid i \geq k,\, \varphi \in
\Gamma\rangle,  \\[.03in]
N_{k}^{-} & := & N_{0,\ldots,k} & := & \langle\varphi^{-1}\rho_{i}\varphi \mid i  \leq k,\,\varphi \in
\Gamma\rangle.
\end{array}\]
These subgroups are the normal closures of $\{\rho_{k},\ldots,\rho_{d-1}\}$ and
$\{\rho_{0},\ldots,\rho_{k}\}$, respectively, in $\Gamma =
\langle\rho_{0},\ldots,\rho_{d-1}\rangle$, and satisfy 
\begin{equation}
\label{nkplus}
\Gamma = N_{k}^{+} \langle\rho_{0},\ldots,\rho_{k-1}\rangle =
\langle\rho_{0},\ldots,\rho_{k-1}\rangle N_{k}^{+},
\end{equation}
\begin{equation}
\label{nkminus}
\Gamma = N_{k}^{-} \langle\rho_{k+1},\ldots,\rho_{d-1}\rangle =
\langle\rho_{k+1},\ldots,\rho_{d-1}\rangle N_{k}^{-}.
\end{equation}

Now let $\mathcal{P}$ be a regular $d$-polytope with automorphism group $\Gamma(\mathcal{P}) =\langle\rho_{0},\ldots,\rho_{d-1}\rangle$, and let $0 \leq k \leq d-1$.   We say that $\mathcal{P}$ has the {\em FAP with respect to its $k$-faces\/} if the above product in (\ref{nkplus}) is semidirect; that is, \[\Gamma(\mathcal{P}) \cong N_{k}^{+} \rtimes \langle\rho_{0},\ldots,\rho_{k-1}\rangle,\]
or, equivalently, the subgroups $N_{k}^{+}$ and $\langle\rho_{0},\ldots,\rho_{k-1}\rangle$ intersect trivially. In this case, if $k=d-1$ we also say that $\mathcal{P}$ has the {\em FAP with respect to its facets\/}.

Dually, $\mathcal{P}$ is said to have the {\em FAP with respect to the co-$k$-faces\/} if the product in (\ref{nkminus}) is semi-direct; that is,
\[\Gamma(\mathcal{P}) \cong N_{k}^{-} \rtimes \langle\rho_{k+1},\ldots,\rho_{d-1}\rangle,\]
or, equivalently, $N_{k}^{-}$ and $\langle\rho_{k+1},\ldots,\rho_{n-1}\rangle$ have trivial intersection. In this case, if $k=0$ or $k=1$, we say that $\mathcal{P}$ has the {\em FAP with respect to its vertex-figures\/} or {\em edge-figures\/}, respectively. Then a regular $d$-polytope $\mathcal{P}$ has the FAP with respect to its co-$k$-faces if and only if its dual $\mathcal{P}^{*}$ has the FAP with respect to its $(d-k-1)$-faces.  (Recall that the dual $\mathcal{P}^{*}$ of $\mathcal{P}$ is the polytope obtained from $\mathcal{P}$ by leaving the face-set of $\mathcal{P}$ unchanged, but reversing the partial order on $\mathcal{P}$.)

We often make use of the fact that if a regular $d$-polytope $\mathcal{P}$ has the
FAP with respect to its $k$-faces, then the $j$-faces of $\mathcal{P}$ with $j>k$ also have the FAP with respect to their $k$-faces, and the co-$j$-faces of $\mathcal{P}$ with $j<k$ have the FAP with respect to their $(k-j-1)$-faces. A similar, dual statement holds for polytopes that have the FAP with respect to their co-$k$-faces.

Lemma~\ref{lemher} below will show that the FAP is also hereditary in another sense. The proof requires the following simple technical lemma.

\begin{lemma}
\label{lemABC}
Let $G$ be a group, and let $A,B,C$ be subgroups of $G$ such that $G=AB$, $A\cap B=C\cap B=1$, $B\unlhd G$, $C\unlhd A$, and $CB\unlhd AB$. Then 
\[ G/CB = AB/CB \cong A/C.\]
\end{lemma}

\begin{proof} 
By the isomorphism theorems for groups,
\[(AB)/B \cong A/(A\cap B) \cong A,\;\; (CB)/B \cong C/(C\cap B) \cong C\] 
and therefore
\[G/CB = AB/CB \cong (AB/B)/(CB/B)\cong A/C.\]
\end{proof}
\smallskip

The first part of the following lemma will be used in the proof of Theorem~\ref{flat44thm}. The second part is the corresponding dual statement.

\begin{lemma} 
\label{lemher}
Let $\mathcal{P}$ be a regular $d$-polytope, $d\geq 4$. \\[.03in]
(a) Let $0\leq k\leq d-3$ and $0\leq l\leq d-k-4$. If $\mathcal{P}$ has the FAP with respect to its co-$k$-faces, and the co-$k$-faces have the $FAP$ with respect to their co-$l$-faces, then $\mathcal{P}$ has the FAP with respect to its co-$(k+l+1)$-faces.\\[.03in]
(b) Let $2\leq l<k\leq d-1$. If $\mathcal{P}$ has the FAP with respect to its $k$-faces, and the $k$-faces have the $FAP$ with respect to their $l$-faces, then $\mathcal{P}$ has the FAP with respect its $l$-faces.
\end{lemma}

\begin{proof}
Let $\Gamma(\mathcal{P}) = \langle\rho_{0},\ldots,\rho_{d-1}\rangle$. As the two parts are related by duality, it suffices to prove the first part. 

The group of the co-$k$-face at the base $k$-face of $\mathcal{P}$ is given by $\langle\rho_{k+1},\ldots,\rho_{d-1}\rangle$. Then our assumptions on $\mathcal{P}$ and its co-$k$-faces imply that 
\[\Gamma(\mathcal{P}) \cong N_{k}^{-} \rtimes \langle\rho_{k+1},\ldots,\rho_{d-1}\rangle\]
and
\[\langle\rho_{k+1},\ldots,\rho_{d-1}\rangle =
\hat{N}_{k+1,\ldots,k+l+1} \rtimes \langle\rho_{k+l+2},\ldots,\rho_{d-1}\rangle,\]
where $\hat{N}_{k+1,\ldots,k+l+1}$ is defined as the normal closure of $\{\rho_{k+1},\ldots,\rho_{k+l+1}\}$ in the subgroup $\langle\rho_{k+1},\ldots,\rho_{d-1}\rangle$. We claim that 
\begin{equation}
\label{herNs}
N_{k+l+1}^{-} = N_{k}^{-}\hat{N}_{k+1,\ldots,k+l+1},
\end{equation}
or equivalently, $N_{k+l+1}^{-}/N_{k}^{-}\cong \hat{N}_{k+1,\ldots,k+l+1}$.
In fact, by the FAP of $\mathcal{P}$ with respect to its co-$k$-faces, $N_k^{-}$ intersects the subgroup $\hat{N}_{k+1,\ldots,k+l+1}$  of $\langle\rho_{k+1},\ldots,\rho_{k+l+1}\rangle$ 
trivially and therefore, 
\[\begin{array}{rcl}
N_{k+l+1}^{-}/ N_{k}^{-}&= &
\langle\rho_{0},\ldots,\rho_{k+l},\varphi\rho_{k+l+1}\varphi^{-1}\!\mid\! \varphi\in\langle\rho_{k+l+2},\ldots,\rho_{d-1}\rangle\rangle/N_{k}^{-}\\
&=&
\langle\rho_{0}N_{k}^{-},\ldots,\rho_{k+l}N_{k}^{-},\varphi\rho_{k+l+1}\varphi^{-1}N_{k}^{-}\mid \varphi\in\langle\rho_{k+l+2},\ldots,\rho_{d-1}\rangle\rangle\\
&=&
\langle\rho_{k+1}N_{k}^{-},\ldots,\rho_{k+l}N_{k}^{-},\varphi\rho_{k+l+1}\varphi^{-1}N_{k}^{-}\mid \varphi\in\langle\rho_{k+l+2},\ldots,\rho_{d-1}\rangle\rangle\\
&=&
\langle\rho_{k+1},\ldots,\rho_{k+l},\varphi\rho_{k+l+1}\varphi^{-1}\mid \varphi\in\langle\rho_{k+l+2},\ldots,\rho_{d-1}\rangle\rangle N_{k}^{-}/N_{k}^{-}\\
&=&\hat{N}_{k+1,\ldots,k+l+1} N_{k}^{-}/N_{k}^{-}\\
&=&\hat{N}_{k+1,\ldots,k+l+1}/(\hat{N}_{k+1,\ldots,k+l+1}\cap N_{k}^{-})\\
&\cong&\hat{N}_{k+1,\ldots,k+l+1}.
\end{array}\]
Thus, $N_{k+l+1}^{-} = N_{k}^{-}\hat{N}_{k+1,\ldots,k+l+1}$.

We now apply Lemma~\ref{lemABC} to complete the proof. More explicitly, since the co-$k$-faces have the FAP with respect to their co-$l$-faces,
\[\begin{array}{rcl}
\Gamma(\mathcal{P})/N_{k+l+1}^{-}
&=&\langle\rho_{k+1},\ldots,\rho_{d-1}\rangle N_{k}^{-} / \hat{N}_{k+1,\ldots,k+l+1}N_{k}^{-} \\
&=&\langle\rho_{k+1},\ldots,\rho_{d-1}\rangle / \hat{N}_{k+1,\ldots,k+l+1}\\
&\cong&\langle\rho_{k+l+2},\ldots,\rho_{d-1}\rangle.
\end{array}\]
Thus $\mathcal{P}$ has the FAP with respect to its co-$(k+l+1)$-faces. 
\end{proof}
\medskip

The flat amalgamation property can be employed to construct flat regular polytopes $\mathcal{P}$ with preassigned faces $\mathcal{P}_{1}$ and co-faces $\mathcal{P}_{2}$.
Suppose $\mathcal{P}_{1}$ is a regular $c$-polytope and $\mathcal{P}_{2}$ a regular $d$-polytope such that the co-$k$-faces $\mathcal{K}$ of $\mathcal{P}_{1}$ are isomorphic to the $(c-k-1)$-faces of $\mathcal{P}_{2}$. Let $\Gamma(\mathcal{P}_{1}) = \langle\alpha_{0},\ldots,\alpha_{c-1}\rangle$, and $\Gamma(\mathcal{P}_{2}) = \langle\beta_{k+1},\ldots,\beta_{k+d}\rangle$, where the subscripts of the standard generators of $\Gamma(\mathcal{P}_{2})$ have been shifted by $k+1$. As in \cite[4F5]{McMSch2002}, consider the subgroup 
$\Gamma := \Gamma(\mathcal{P}_{1},\mathcal{P}_{2}) := \langle\rho_0,\ldots,\rho_{k+d}\rangle$ of the direct product $\Gamma(\mathcal{P}_{1}) \times \Gamma(\mathcal{P}_{2})$ where the generators $\rho_i$ are defined by
\begin{equation}
\label{flatamalgen}
\rho_{i} := \left\{  
\begin{array}{lcl}
(\alpha_{i},1) & \mbox{if} & 0 \leq i \leq k,  \\
(\alpha_{i},\beta_{i}) & \mbox{if} & k+1 \leq i \leq c-1,  \\
(1,\beta_{i}) & \mbox{if} & c \leq i \leq k+d.
\end{array}  \right.
\end{equation}
Then $\Gamma$ is called the $(k+1)$-{\it mix\/} of the groups $\Gamma(\mathcal{P}_{1})$ and $\Gamma(\mathcal{P}_{2})$, denoted $\Gamma(\mathcal{P}_{1})\!\diamond_{k+1}\!\Gamma(\mathcal{P}_{2})$. If $\Gamma$ happens to be a C-group (as in the situation described below), then the corresponding regular polytope is called the $(k+1)$-mix of the polytopes $\mathcal{P}_1$ and $\mathcal{P}_2$ denoted $\mathcal{P}_{1}\!\diamond_{k+1}\!\mathcal{P}_{2}$.
The specific nature of the generators immediately shows that $N_{k}^{-}$ and $N_{c}^{+}$ lie in $\Gamma(\mathcal{P}_{1})\times\{1\}$ or $\{1\}\times\Gamma(\mathcal{P}_{2})$, respectively, and thus commute at the level of elements.

The following theorem summarizes the properties of this mix construction (see \cite[Thm. 4F9]{McMSch2002}). 

\begin{theorem}
\label{amalprops}
Let $c,d \geq 2$, and let $0 \leq k \leq c-2$ and $k \geq c-d$.  Let $\mathcal{P}_{1}$
be a regular $c$-polytope with group $\Gamma(\mathcal{P}_{1}) = \langle\alpha_{0},\ldots,\alpha_{c-1}\rangle$, and let $\mathcal{P}_{2}$ be a regular $d$-polytope with $\Gamma(\mathcal{P}_{2}) = \langle\beta_{k+1},\ldots,\beta_{k+d}\rangle$ (with subscripts shifted by $k+1$). Suppose that the co-$k$-faces $\mathcal{K}$ of $\mathcal{P}_{1}$ are isomorphic to the $(c-k-1)$-faces of $\mathcal{P}_{2}$, and that $\mathcal{P}_{1}$ has the FAP with respect to its co-$k$-faces and $\mathcal{P}_{2}$ has the FAP with respect to its $(c-k-1)$-faces. Then the $(k+1)$-mix $\Gamma = \Gamma(\mathcal{P}_{1},\mathcal{P}_{2})$ of $\Gamma(\mathcal{P}_{1})$ and $\Gamma(\mathcal{P}_{2})$ is a string C-group of rank $k+d+1$. Let $\mathcal{P} := \mathcal{P}_{1} \diamond_{k+1} \mathcal{P}_{2}$ be the corresponding regular $(k+d+1)$-polytope with group $\Gamma(\mathcal{P}) = \Gamma$. Then $\mathcal{P}$ has the following properties:\\[.02in]
(a)\, The $c$-faces of $\mathcal{P}$ are isomorphic to $\mathcal{P}_{1}$, and the
co-$k$-faces are isomorphic to $\mathcal{P}_{2}$.  More generally, if $r \leq k$, \ 
$s \geq c$, \ $s-r \geq 4$, and $\{F_{-1},F_0,\ldots,F_{k+d+1}\}$ denotes the base flag of $\mathcal{P}$, then the section $F_{s}/F_{r}$ of $\mathcal{P}$ is
isomorphic to the regular $(s-r-1)$-polytope $\mathcal{Q} := \mathcal{Q}_{1} \diamond_{k-r}\mathcal{Q}_{2}$ which is constructed from the co-$r$-face $\mathcal{Q}_{1}$ of $\mathcal{P}_{1}$ and $(s-k-1)$-face $\mathcal{Q}_{2}$ of $\mathcal{P}_{2}$ in the same way as $\mathcal{P}$ is from $\mathcal{P}_{1}$ and $\mathcal{P}_{2}$.\\[.02in]
(b)\, $\mathcal{P}$ has the FAP with respect to both its co-$k$-faces and its
$c$-faces.  In particular,
\[  \Gamma(\mathcal{P}) \cong 
N_{k}^{-} \rtimes \Gamma(\mathcal{P}_{2}) \cong
N_{c}^{+} \rtimes \Gamma(\mathcal{P}_{1}) \cong 
(N_{k}^{-} \times N_{c}^{+}) \rtimes \Gamma(\mathcal{K}),  \]
where $N_{k}^{-}$ and $N_{c}^{+}$ are the normal closures of 
$\{\alpha_{0},\ldots,\alpha_{k}\}$ in $\Gamma(\mathcal{P}_{1})$ and 
$\{\beta_{c},\ldots,\beta_{k+d}\}$ in $\Gamma(\mathcal{P}_{2})$, respectively.\\[.02in]
(c)\, $\mathcal{P}$ is $(k,c)$-flat.
\end{theorem}

The toroidal regular polyhedra $\{4,4\}_{(s,t)}$, with $s\geq 2$, $t=0$ or $s=t\geq 2$, have been extensively studied (see Coxeter \& Moser~\cite[Ch. 8]{CoxMos1980}). 
They can be derived from the square tessellation of the plane with vertex set $\mathbb{Z}^2$ as the quotient by the lattices spanned by the vectors $(s,t)$, $(-t,s)$, respectively. Here we briefly discuss the polyhedra $\{4,4\}_{(s,t)}$ whose automorphism groups are 2-groups.

\begin{lemma}
\label{lem44}
The regular polyhedra of type $\{4,4\}$ with 2-groups as automorphism groups are the polyhedra $\{4,4\}_{(2^e,0)}$ and $\{4,4\}_{(2^e,2^e)}$, $e\geq 1$, with group orders $2^{2e+3}$ and $2^{2e+4}$, respectively. 
\end{lemma}

\begin{proof}
A regular polyhedron $\{4,4\}_{(s,t)}$, with $s\geq 2$, $t=0$ or $s=t\geq 2$, has an automorphism group of order $8(s^2+t^2)$. For this group to be 2-group, $s^2+t^2$ must be a power of 2. Using simple arithmetic modulo 4, it is not difficult to see that each power of 2 can be written in exactly one way as a sum of two squares of non-negative numbers, up to interchanging the two summands:\ if $k=2e$ is even, then $2^k = (2^e)^2 +0^2$, and if $k=2e+1$ is odd, then $2^k = (2^e)^2 +(2^e)^2$. Thus, if the order of the group is $2^n$, $n\geq 4$, the equation $2^{n-3}=s^2+t^2$ gives the parameters $(s,t)=(2^e,0),(2^e,2^e)$ with $n=2e+3,\,2e+4$, respectively. This proves the lemma.
\end{proof}

For $n\geq 5$, we let $\{4,4\}^{(n)}$ denote the unique regular polyhedron of type $\{4,4\}$ with an automorphism group of order $2^n$. Thus, 
\begin{equation} 
\{4,4\}^{(n)}=\{4,4\}_{(2^e,0)}\mbox{ or }\{4,4\}_{(2^e,2^e)},
\end{equation} 
according as $n$ is odd, $n-3=2e$, or $n$ is even, $n-4=2e$. We similarly use the notation $[4,4]^{(n)}$, $[4,4]_{(2^e,0)}$ or $[4,4]_{(2^e,2^e)}$ for the respective automorphism groups. Each regular polyhedra $\{4,4\}^{(n)}$ is self-dual and has a polarity $\omega$ (duality of order 2) that fixes the base flag. If we let $[4,4]^{(n)} = \langle\alpha_0,\alpha_1,\alpha_2\rangle$, then conjugation by $\omega$ in the extended group (of all automorphisms and dualities) swaps the generators $\alpha_0$ and $\alpha_2$ while keeping the generator $\alpha_1$ fixed.

The polyhedra $\{4,4\}^{(n)}$ all have the flat amalgamation property with respect to both, the vertex-figures and the facets~\cite{Sch1988}. In particular, if $N(\alpha_0)$ denotes the normal closure of $\alpha_0$ in $[4,4]^{(n)}$, then $N(\alpha_0)$ has order $2^{n-3}$ and
\begin{equation}
\label{torfap1}
[4,4]^{(n)} = N(\alpha_0)\rtimes\langle\alpha_1,\alpha_2\rangle 
\cong N(\alpha_0) \rtimes D_4 .
\end{equation}
Similarly, 
\begin{equation}
\label{torfap2}
[4,4]^{(n)} = N(\alpha_2)\rtimes\langle\alpha_0,\alpha_1\rangle 
\cong N(\alpha_2)\rtimes D_4 ,
\end{equation}
where $N(\alpha_2)$ denotes the normal closure of $\alpha_2$ in $[4,4]^{(n)}$.
\medskip

\section{Regular polytopes with 2-groups}
\label{regpols2}
\medskip

In this section, we construct flat regular $d$-polytopes $\mathcal P$ whose automorphism groups are finite 2-groups and whose successive sections of rank 3 as in (\ref{rk3sections}) are toroidal maps of the kind described in the previous section. Any such polytope $\mathcal P$ can be assigned the generalized Schl\"afli symbol
\begin{equation}
\label{all4stypes}
\{\{4,4\}^{(n_3)},\{4,4\}^{(n_4)},\ldots,\{4,4\}^{(n_{d})}\},
\end{equation}
where here the toroidal polyhedron $\{4,4\}^{(n_j)}$ (whose automorphism group has order $2^{n_j}$) gives the isomorphism type of the section 
$F_{j}/F_{j-4}$ of $\mathcal P$ for $j=3,\ldots,d$.  We will show that any sequence of numbers $n_3,\ldots,n_d$, with $n_{j}\geq 5$ for $j=3,\ldots,d$, can be preassigned. 

\section{Counted theorem and complete author proof}
\begin{theorem}
\label{flat44thm}
Let $d\geq 3$, and let $n_3,\ldots,n_d\geq 5$. Then there exists a finite regular $d$-polytope $\mathcal P$ with the following properties:\\[.02in]
(a) $\mathcal P$ is of type $\{\{4,4\}^{(n_3)},\{4,4\}^{(n_4)},\ldots,\{4,4\}^{(n_{d})}\}$.\\[.02in]
(b) The automorphism group of $\mathcal P$ is a 2-group of order 
$2^{n_3+n_4+\ldots+n_d - 3(d-3)}$.\\[.02in]
(c) $\mathcal P$ is flat if $d\geq 4$; in fact, $\mathcal P$ is $(d-4,d-1)$-flat if $d\geq 4$.\\[.02in]
(d) $\mathcal P$ has the FAP with respect to its co-$(d-4)$-faces, its co-$(d-3)$-faces, and its facets.
\end{theorem}

\begin{proof}
We use Theorem~\ref{amalprops} and exploit the fact that the toroidal maps $\{4,4\}^{(n_j)}$ all have the flat amalgamation property with respect to both, the facets and the vertex-figures. 

Our proof is by induction on the rank $d$. For each dimension $d\geq 3$, we construct a regular polytope ${\mathcal P}_d$ with the desired properties. For $d=3$ there is nothing to prove: ${\mathcal P}_{3}=\{4,4\}^{(n_3)}$. For $d\geq 4$, the polytope ${\mathcal P}_{d}$ will be derived by amalgamating ${\mathcal P}_{d-1}$ and $\{4,4\}^{(n_d)}$ along the $2$-sections which are the co-$(d-4)$-face of ${\mathcal P}_{d-1}$ and the $2$-face of $\{4,4\}^{(n_d)}$, respectively. 

It is instructive to first consider the case $d=4$. Set ${\mathcal Q}_1:=\{4,4\}^{(n_3)}$ and 
${\mathcal Q}_2:=\{4,4\}^{(n_4)}$. Then ${\mathcal Q}_1$ has the FAP with respect to its vertex-figures, and ${\mathcal Q}_2$ has the FAP with respect to its 2-faces. In particular, if we let $\Gamma(\mathcal{Q}_1)=\langle\alpha_0,\alpha_1,\alpha_2\rangle$ and $\Gamma(\mathcal{Q}_2)=\langle\beta_1,\beta_2,\beta_3\rangle$, then 
\[\Gamma(\mathcal{Q}_1)=N(\alpha_0)\rtimes\langle\alpha_1,\alpha_2\rangle,\quad
\Gamma(\mathcal{Q}_2)=N(\beta_3)\rtimes\langle\beta_1,\beta_2\rangle,\]
where $\langle\alpha_1,\alpha_2\rangle\cong\langle\beta_1,\beta_2\rangle\cong D_4$ and $N(\alpha_0)$ and $N(\beta_3)$ are the normal closures of $\alpha_0$ in 
$\Gamma(\mathcal{Q}_1)$ and $\beta_3$ in $\Gamma(\mathcal{Q}_2)$, respectively. Here, the subgroups $N(\alpha_0)$ and $N(\beta_3)$ have orders $2^{n_{3}-3}$ and $2^{n_{4}-3}$, respectively. Now, by Theorem~\ref{amalprops} above (with $k=0$), there exists a $(0,3)$-flat regular 4-polytope, denoted $\mathcal{P}_4$,  such that the facets and vertex-figures of ${\mathcal P}_4$ are isomorphic to ${\mathcal Q}_1$ and ${\mathcal Q}_2$, respectively. Thus ${\mathcal P}_4$ is of the desired type $\{\{4,4\}^{(n_3)},\{4,4\}^{(n_4)}\}$. If we let $\Gamma({\mathcal P}_4)=\langle\rho_0,\ldots,\rho_3\rangle$, with generators $\rho_i$ as in (\ref{flatamalgen}), and let $N_{0}\,(\cong N(\alpha_0))$ and $N_{3}\,(\cong N(\beta_3))$ denote the normal closures of $\rho_0$ and $\rho_3$ in the facet subgroup $\langle\rho_0,\rho_1,\rho_2\rangle$ and vertex-figure subgroup $\langle\rho_1,\rho_2,\rho_3\rangle$, respectively, then $N_{0}$ and $N_{3}$ are also the normal closures of $\rho_0$ and $\rho_3$ in the full group $\Gamma({\mathcal P}_4)$ and
\[\Gamma({\mathcal P}_4)
\cong N_0\rtimes\Gamma(\mathcal{Q}_2)
\cong N_3\rtimes\Gamma(\mathcal{Q}_1)
\cong (N_0\times N_3)\rtimes D_4 ,\]
where the dihedral group $D_4$ is given by $\langle\rho_1,\rho_2\rangle$. Moreover, ${\mathcal P}_4$ has the FAP with respect to its vertex-figures and its facets. Note that $\Gamma({\mathcal P}_4)$ has order
\[2^{(n_{3}-3)+(n_{4}-3)+3}=2^{n_3+n_4-3}=2^{n_3+n_4-3(d-3)}.\]

In order for us to be able to proceed to the next rank via the outlined amalgamation process, we must verify that $\mathcal{P}_4$ also has the FAP with respect to its edge-figures. But this now follows from Lemma~\ref{lemher}(a). In fact, since $\mathcal{P}_4$ has the FAP with respect to its vertex-figures, and the vertex-figures of $\mathcal{P}_4$ have the FAP with respect to their vertex-figures, Lemma~\ref{lemher}(a) shows that $\mathcal{P}_4$ also has the FAP with respect to its edge-figures. In particular, 
\[\Gamma({\mathcal P}_4) = N_{0,1} \langle\rho_2,\rho_3\rangle 
\cong N_{0,1}\rtimes \langle\rho_2,\rho_3\rangle,\]
where $N_{0,1}$ is the normal closure of $\{\rho_0,\rho_1\}$ in $\Gamma({\mathcal P}_4)$. In this context, (\ref{herNs}) takes the form 
\begin{equation}
\label{n01n1hat}
N_{0,1}=N_0 \,\hat{N}_{1} \cong N_0 \rtimes \hat{N}_{1},
\end{equation}
where $\hat{N}_{1}$ is the normal closure of $\rho_1$ in the vertex-figure subgroup $\langle\rho_1,\rho_2,\rho_3\rangle$.

Note that the semidirect product in (\ref{n01n1hat}) is not a direct product; the elements $\rho_{1}\rho_{0}\rho_{1}$ of $N_0$ and $\rho_{2}\rho_{1}\rho_{2}$ of $\hat{N}_{1}$ do not commute.
In fact, if we assume the contrary and keep in mind that $\rho_1\rho_2$ has order 4, then the equation
$\rho_{2}\rho_{1}\rho_{2}\cdot\rho_{1}\rho_{0}\rho_{1} \cdot\rho_{2}\rho_{1}\rho_{2}
=\rho_{1}\rho_{0}\rho_{1}$ simplifies to $\rho_{2}\rho_{1}\rho_{2}\rho_{0}\rho_{2}\rho_{1}\rho_{2}=\rho_0$ and therefore $(\rho_1\rho_0)^2=1$, contradicting the order of $\rho_0 \rho_1$. %This forces $\mathcal{Q}_1$ to be $\{4,4\}_{(1,0)}$ or $\{4,4\}_{(1,1)}$. However, these regular maps are not polyhedra and are excluded as possible sections.

For the inductive step of the construction, let $d\geq 5$ and suppose there exists a regular $(d-1)$-polytope ${\mathcal P}_{d-1}$ with the properties listed in the theorem. Set ${\mathcal Q}_1:=\mathcal{P}_{d-1}$ and ${\mathcal Q}_2:=\{4,4\}^{(n_d)}$. Then ${\mathcal Q}_1$ has the FAP with respect to its co-$(d-4)$-face, and ${\mathcal Q}_2$ has the FAP with respect to its $2$-faces. Now, by Theorem~\ref{amalprops} above, there exists a $(d-4,d-1)$-flat regular $d$-polytope, denoted $\mathcal{P}_d$,  such that the facets and co-$(d-4)$-faces of ${\mathcal P}_d$ are isomorphic to ${\mathcal Q}_1$ and ${\mathcal Q}_2$, respectively. If we let $\Gamma({\mathcal P}_d)=\langle\rho_0,\ldots,\rho_{d-1}\rangle$, then 
\[\Gamma({\mathcal P}_d)
\cong N_{0,1,\ldots,d-4}\rtimes\Gamma(\mathcal{Q}_2)
\cong N_{d-1}\rtimes\Gamma(\mathcal{Q}_1)
\cong (N_{0,1,\ldots,d-4}\times N_{d-1})\rtimes D_4 ,\]
where here $N_{0,1,\ldots,d-4}$ denotes the normal closure of $\{\rho_0,\rho_1,\ldots,\rho_{d-4}\}$ in the facet subgroup $\langle\rho_0,\ldots,\rho_{d-2}\rangle$, and $N_{d-1}$ is the normal closure of $\rho_{d-1}$ in the co-$(d-4)$-face subgroup $\langle\rho_{d-3},\rho_{d-2},\rho_{d-1}\rangle$. These subgroups $N_{0,1,\ldots,d-4}$ and $N_{d-1}$ are also the normal closures of $\{\rho_0,\rho_1,\ldots,\rho_{d-4}\}$ and $\rho_{d-1}$ in the full group $\Gamma({\mathcal P}_d)$, respectively. As $\mathcal{Q}_{1}=\mathcal{P}_{d-1}$, the order of the automorphism group is given by 
\[|\Gamma({\mathcal P}_d)| = |N_{d-1}|\cdot |\Gamma(\mathcal{P}_{d-1})|
=2^{n_{d}-3}\cdot 2^{n_3+\ldots +n_{d-1} -3(d-4)} = 2^{n_3+\ldots +n_{d} -3(d-3)}.\]
By Theorem~\ref{amalprops}, ${\mathcal P}_d$ also has the FAP with respect to both, its co-$(d-4)$-faces and its facets.

It remains to prove that ${\mathcal P}_d$ also has the FAP with respect to its co-$(d-3)$-faces. We can proceed as for rank 4 and appeal to Lemma~\ref{lemher}. Then it follows that $\mathcal{P}_d$ has the FAP with respect to its co-$(d-3)$-faces and 
\[\Gamma({\mathcal P}_d) = N_{0,\ldots,d-3}\, \langle\rho_{d-2},\rho_{d-1}\rangle 
\cong N_{0,\ldots,d-3}\rtimes \langle\rho_{d-2},\rho_{d-1}\rangle,\]
where $N_{0,\ldots,d-3}$ is the normal closure of $\{\rho_0,\ldots,\rho_{d-3}\}$ in $\Gamma({\mathcal P}_d)$. Here, (\ref{herNs}) takes the form 
\begin{equation}
\label{theNs}
N_{0,\ldots,d-3}=N_{0,\ldots,d-4} \,\hat{N}_{d-3}\cong N_{0,\ldots,d-4} \rtimes \hat{N}_{d-3},
\end{equation}
where $\hat{N}_{d-3}$ is the normal closure of $\rho_{d-3}$ in the co-$(d-4)$-face subgroup $\langle\rho_{d-3},\rho_{d-2},\rho_{d-1}\rangle$.
Note that the semidirect product in (\ref{theNs}) is not direct, as the elements $\rho_{d-3}\rho_{d-4}\rho_{d-3}$ of $N_{0,\ldots,d-4}$ and $\rho_{d-2}\rho_{d-3}\rho_{d-2}$ of $\hat{N}_{d-3}$ do not commute, for similar reasons as before.

Thus the regular $d$-polytope ${\mathcal P}_d$ has the desired properties, and the inductive proof of the theorem is complete.
\end{proof}

\begin{thebibliography}{99}
\bibitem[CM80]{CoxMos1980} H.~S.~M.~Coxeter and W.~O.~J.~Moser, \textit{Generators and Relations for Discrete Groups}, 4th Edition, Springer, 1980.
\bibitem[MS02]{McMSch2002} P.~McMullen and E.~Schulte, \textit{Abstract Regular Polytopes}, Cambridge University Press, 2002.
\bibitem[Sch88]{Sch1988} E.~Schulte, Amalgamation of regular incidence-polytopes, \textit{Proc. London Math. Soc.\/} (3) {\bf 56} (1988), 303--328.
\end{thebibliography}
\end{document}
